\section{P versus NP: rank certificate payment}\label{sec:pnp}

For the twins (\cref{sec:reduction,sec:firstcause}) and for the Riemann hypothesis
(\cref{sec:riemann}) the green machine carried us to a node that we then accepted by decree ---
the first cause \axiom{}. Here, uniquely among the branches, the machine behaves differently:
the reading of $\mathrm{P}$ versus $\mathrm{NP}$~\cite{clay-pnp} in the language of rank is a
\emph{green theorem}, and no decree is needed at all.

The reading is physical. An $\mathrm{NP}$ problem is the \emph{full payment} of all rank
certificates: to solve the search is to account for every witness individually. $\mathrm{P}$ is a
\emph{fast traversal} of the rank: the motion of an engine that flies past the rank without
visiting all its states. The inequality between them is that a fast traversal physically cannot
pay for an unbounded family of certificates --- pure pigeonhole. All of this is machine-checked
without any axiom. We emphasise at once what is \emph{not} claimed: this is not classical
$\mathrm{P}\neq\mathrm{NP}$. The abstract layer of complexity classes is plastic (it can be tailored
to coincide or to separate for free), so nothing about the genuine classical classes follows from
it, and we make no such claim.

\subsection*{The two sides in the language of rank}

\begin{definition}[\code{RankFastTraversal} / \code{FullRankCertificatePayment}]\label{def:pnp-sides}
The \emph{P side} is \code{RankFastTraversal}: in a ranked graph every legal path \code{PathN} of
length $n$ from $x$ satisfies $n \le \operatorname{lexRank}(x)$. The \emph{NP side} is
\code{FullRankCertificatePayment}: the existence of an injective finite-key compression of the
\emph{entire} family of genealogies, each certificate accounted for individually. An
\code{UnboundedCertificateSupply} is a family whose index set is infinite.
\end{definition}

The P side is a gift of the construction, not a hypothesis.

\begin{theorem}[\code{rankFastTraversal_holds}]\label{thm:rankFastTraversal_holds}
\green{}\ Every ranked graph traverses the rank fast: for any ranked graph $G$, every legal path
\code{PathN} of length $n$ from $x$ satisfies $n \le \operatorname{lexRank}(x)$.
\end{theorem}

\emph{Proof idea.} A direct consequence of the fact that the length of a legal path does not
exceed the starting rank (\code{len_le_lexRank}, \cref{sec:engine}); verification is linear in the
rank. Verifying a single certificate is likewise always cheap (\code{verificationEasy_always},
\green{}); what is expensive is accounting for all of them at once.

\subsection*{The heart of the inequality: pigeonhole}

\begin{theorem}[\code{no_fullPayment_of_unboundedSupply}]\label{thm:no_fullPayment_of_unboundedSupply}
\green{}\ A fast finite-key engine cannot fully pay for an unbounded family of certificates: if a
genealogy family $F$ satisfies \code{UnboundedCertificateSupply} (its index set is infinite), then
\code{FullRankCertificatePayment} for $F$ fails --- no injective finite-key compression of $F$
exists.
\end{theorem}

\emph{Proof idea.} Pure pigeonhole. A finite key compressing an infinite supply is forced to
collide two distinct witnesses into one; otherwise it would distinguish infinitely many objects
with finitely many labels. A collision is double bookkeeping, an impossible payment. It is cheap to
traverse the rank and expensive to touch every state, and the second is incompatible with the
first under an infinite supply. This is the green theorem that makes everything work, and it owes
nothing to \axiom{}.

It remains to exhibit a scale where the supply is genuinely infinite --- and there is one, without
any hypothesis about the twins.

\begin{theorem}[\code{concreteSupply_unbounded_smallScale}]\label{thm:concreteSupply_unbounded_smallScale}
\green{}\ For $A \le 4$ the concrete family of certificates is infinite: the index type of
\code{concreteFamily}\,$A\,1$ is infinite, witnessed by the injection of the pentadic chain
\code{fiveAdicChainFlow} --- the same chain that refuted the branch $A \le 4$ in the reduction.
\end{theorem}

\begin{corollary}[\code{concrete_noFullPayment_smallScale}]\label{thm:concrete_noFullPayment_smallScale}
\green{}\ For $A \le 4$ full payment of the concrete family is impossible:
\code{FullRankCertificatePayment}\,(\code{concreteFamily}\,$A\,1$) fails.
\end{corollary}

\begin{proof}
Apply \cref{thm:no_fullPayment_of_unboundedSupply} to the infinite supply of
\cref{thm:concreteSupply_unbounded_smallScale}.
\end{proof}

\subsection*{``NP $=$ full payment'' as a theorem}

To make the title a theorem rather than an image, we tie ``full payment'' to local P-success --- at
every scale.

\begin{theorem}[\code{concrete_localPSuccess_iff_fullPayment}]\label{thm:concrete_localPSuccess_iff_fullPayment}
\green{}\ Local P-success is equivalent to full payment of certificates: for all $A, M_0$,
\[
   \code{LocalPSuccess}\,(\code{concreteProblem}\,A\,M_0)
   \iff
   \code{FullRankCertificatePayment}\,(\code{concreteFamily}\,A\,M_0).
\]
\end{theorem}

\emph{Proof idea.} Both alternatives for resolving a collision --- a legal cycle and an impossible
payment --- have already been burnt green (\code{no_extendedFlowResolutionAlternative}): the
resolver cannot cheat, and the only way to resolve collisions is to account for each one honestly.
The full inequality now assembles into a single green statement.

\begin{theorem}[\code{pnp_rank_separation_smallScale}]\label{thm:pnp_rank_separation_smallScale}
\green{}\ For $A \le 4$ five statements hold simultaneously:
\code{RankFastTraversal}\,(\code{concreteGraph}\,$A\,1$) (the engine traverses the rank fast);
\code{VerificationEasy} for every certificate (checking is easy);
\code{UnboundedCertificateSupply}\,(\code{concreteFamily}\,$A\,1$) (the supply is unbounded);
$\lnot\,\code{FullRankCertificatePayment}\,(\code{concreteFamily}\,A\,1)$ (full payment is
impossible); and \code{LocalSearchIncompressible}\,(\code{concreteProblem}\,$A\,1$) (local search
is incompressible).
\end{theorem}

The five facets --- fast traversal, easy verification, infinite supply, impossibility of full
payment, incompressibility --- converge in one green statement. This is ``$\mathrm{P}\neq\mathrm{NP}$''
\emph{in the rank model}, proved without axioms.

\paragraph{The split across scales.}
This does not contradict the twins decree, which on its own scale precisely \emph{resolves}
collisions. Local P-success and incompressibility are strict negations, but they live on different
scales: the decree works for $A \ge 5$, where on its own scale the first cause fully pays the
certificates (\code{decreedScale_fullPayment}, \yellow{}), while incompressibility lives at
$A \le 4$, where the supply is infinite and payment impossible --- unconditionally green, owing
nothing to \axiom{}. Note the asymmetry with Riemann: there the boundary was \emph{needed} for the
carrying lemma; here the separation does not need it. We leave the asymmetry visible rather than
insert an unused hypothesis for symmetry's sake.

\subsection*{Why there is no third boundary of the decree}

For the twins and Riemann the decree was an honest second boundary. It is natural to ask for a
third, a $\mathrm{P}/\mathrm{NP}$ boundary; the answer is no, and this is proved machine-wise. All
three conceivable forms of such a field are green verdicts: the universal form (in every good-faith
frame a P-solver extracts a resolver) is \emph{refutable} (\code{pnpLawUniversal_refuted}); the
decider form (reconstructing a resolver from a bare decider) is likewise \emph{refutable}
(\code{pnpLawDeciderGated_refuted}), since a bare \code{PDecider} exists classically for any
language and the resolver extracted from it crashes against incompressibility at $A \le 4$; and the
existential form is \emph{already proved} (\code{pnpLawExistential_green}), so decreeing it would be
empty. The cost mirror collapses on its own.

\begin{proposition}[\code{pnpLaw_asserts_separation}]\label{thm:pnpLaw_asserts_separation}
\green{}\ The $\mathrm{P}/\mathrm{NP}$ law is equivalent to the separation \emph{without any
boundary}, and both sides are already theorems; hence no honest third field of the decree exists.
\end{proposition}

\emph{Proof idea.} The abstract layer of complexity classes is plastic: it can be tailored so the
classes coincide for free (\code{allPFrame}) and so they separate for free (\code{constantsFrame}).
Along the way an adversarial audit uncovered a fourth episode of vacuity: a bare \code{PDecider}
carries no complexity content, so over the incompressible node the decider-gated fronts
\code{FaithfulSelfReductionFront} and \code{CurrentExtractionFront} are classically empty and their
separation conclusions are \warn{}~vacuous. Only the decider channel is affected; the \code{InP}
bridge \code{Step00ToClassicalBridge} remains an honest conditional form. The genuinely missing
object is not a proposition over abstract frames but a data-anchored real machine model --- the same
lesson taught by the spectral audit of Riemann.

\subsection*{The epistemic complement}

There is an epistemic twin of this prohibition, mirroring the unknowability results for Collatz and
Yang--Mills (\cref{sec:collatz}).

\begin{theorem}[\code{pnpCause_unknowable}]\label{thm:pnpCause_unknowable}
\green{}\ An internal, finite-fuel resolution of $\mathrm{P}/\mathrm{NP}$ would be a perpetual
engine: it is unreachable from inside the system, beyond the same horizon as the cause of Collatz.
\end{theorem}

\emph{Proof idea.} Green, without a decree; the axiom-taint does not grow. ``It cannot be known
from inside'' is here a theorem, not a slogan: from within a finite-fuel model the resolution
cannot be self-grounded, since the attempt would cost exactly a perpetual engine
(\cref{sec:engine}).

\paragraph{Honest status.}
Among the branches, this is the single case where the inequality is a green theorem \emph{without}
any decree: in the rank model the fast engine provably does not pay for the unbounded supply. The
thermodynamic intuition --- full accounting costs more than fast motion --- is proved in the rank
world; its transfer to the world of Turing machines remains the missing data anchor. Classical
$\mathrm{P}\neq\mathrm{NP}$~\cite{clay-pnp} is neither proved nor claimed.
